%=====================================================================
\section{Geometry of the Complex Simplex Network}
\label{ch:simplex_geometry}
%=====================================================================

This chapter develops the geometry of the simplex network using only combinatorial and algebraic structures. No physical terminology is used. Physical identifications (which vertex type is ``spatial,'' which force is ``strong,'' etc.) are deferred to Chapter~\ref{ch:geometry}.

\subsection{The simplex}

A 4-simplex $\sigma$ has 5 vertices. Each vertex carries one complex number. The state of a simplex is
\begin{equation}
z = (z_0, z_1, z_2, z_3, z_4) \in \mathbb{C}^5.
\end{equation}
After projectivization ($z \sim \lambda z$ for $\lambda \in \mathbb{C}^*$), this is a point in $\mathbb{CP}^4$.

\subsection{The network}

$N$ such simplices form a fully connected weighted network. The state of the network is an $N \times 5$ matrix $\Psi$, where row $i$ is the $\mathbb{C}^5$ vector of simplex $i$. The Gram matrix
\begin{equation}
G_{ij} = \langle\psi_i|\psi_j\rangle = \sum_{k=0}^{4} \bar{z}_{i,k}\, z_{j,k}
\end{equation}
is $N \times N$, Hermitian, positive semidefinite, and has $\text{rank}(G) \leq 5$.

Eigenvalue decomposition of $G$ yields at most 5 significant eigenvalues. The remaining $N - 5$ eigenvalues are $0^+$ (ghosts) or exactly 0. The effective dimension of the network is 5.

\subsection{The (3,2) vertex labeling}

Label 3 vertices as \textbf{A-type} and 2 as \textbf{B-type}: $(A_1, A_2, A_3, B_1, B_2)$. This labeling is the unique non-trivial atomic decomposition of $d = 5$, proven in Chapter~\ref{ch:whyd5}: the only pair of non-repeating (atomic) dimensions summing to 5 is $(3, 2)$.

\begin{definition}[Notation]
Throughout this chapter: $n_A = 3$, $n_B = 2$, $d = n_A + n_B = 5$.
\end{definition}

\subsection{Faces}

A \textbf{face} of the simplex is a subset of 4 vertices (a tetrahedron). Two simplices are \emph{adjacent} when they share a face: 4 vertices in common, 1 different.

Each simplex has $\binom{5}{4} = 5$ faces. With the $(3,2)$ labeling, there are exactly two face types:

\begin{center}
\begin{tabular}{clccl}
\hline
Type & Composition & Count & Changed vertex & Neighbor direction \\
\hline
AAAB & $\{A_1, A_2, A_3, B_j\}$ & $\binom{3}{3}\binom{2}{1} = 2$ & a B vertex & B-direction \\
AABB & $\{A_i, A_j, B_1, B_2\}$ & $\binom{3}{2}\binom{2}{2} = 3$ & an A vertex & A-direction \\
\hline
\multicolumn{2}{r}{Total} & 5 & & \\
\hline
\end{tabular}
\end{center}

\begin{theorem}[Neighbor structure]
A simplex has 2 B-direction neighbors and 3 A-direction neighbors. The adjacency structure is $2 + 3 = 5$.
\end{theorem}

\begin{remark}
The numbers 2 and 3 are not chosen but derived: they are the dimensions of the only non-trivial atomic pair summing to 5.
\end{remark}

\subsection{Hinges}

\subsubsection{Definition}

Curvature in a simplicial complex lives at codimension-2 surfaces. In 4 dimensions, codimension-2 = dimension 2 = \textbf{triangle}. A triangle formed by 3 vertices of the simplex is called a \textbf{hinge}.

The hinge area is
\begin{equation}
A_h = \sqrt{\det(G_h)},
\end{equation}
where $G_h$ is the $3 \times 3$ Gram sub-matrix of the 3 hinge vertices.

\begin{theorem}[1 hinge = 1 bit]
\label{thm:hinge_bit}
Each hinge carries exactly 1 bit of geometric information: $\det(G_h) > 0$ (non-degenerate, ``exists'') or $\det(G_h) = 0$ (degenerate).
\end{theorem}

\subsubsection{Hinges per simplex}

A simplex has $\binom{5}{3} = 10$ hinges. With the $(3,2)$ labeling:

\begin{center}
\begin{tabular}{clc}
\hline
Type & Vertices & Count \\
\hline
AAA & 3 A-vertices & $\binom{3}{3}\binom{2}{0} = 1$ \\
AAB & 2 A + 1 B & $\binom{3}{2}\binom{2}{1} = 6$ \\
ABB & 1 A + 2 B & $\binom{3}{1}\binom{2}{2} = 3$ \\
BBB & 3 B-vertices & $\binom{2}{3} = 0$ \\
\hline
\multicolumn{2}{r}{Total} & 10 \\
\hline
\end{tabular}
\end{center}

\begin{theorem}[BBB impossibility]
\label{thm:bbb}
The BBB hinge does not exist: $\binom{n_B}{3} = \binom{2}{3} = 0$. There are only $n_B = 2$ B-vertices, but a triangle requires 3.
\end{theorem}

\subsubsection{Hinges per face}

Each face (4 vertices) contains $\binom{4}{3} = 4$ hinges. The hinge content depends on face type:

\begin{center}
\begin{tabular}{ccccc}
\hline
Face type & AAA & AAB & ABB & Total \\
\hline
AAAB & 1 & 3 & 0 & 4 \\
AABB & 0 & 2 & 2 & 4 \\
\hline
\end{tabular}
\end{center}

\begin{theorem}[Hinge-face duality]
\label{thm:hinge_face}
Each of the 10 hinges belongs to exactly 2 of the 5 faces: $5 \times 4 / 2 = 10$.
\end{theorem}

\begin{corollary}[Information content]
1 hinge = 1 bit. 1 face = 4 bits. 1 simplex surface = 10 bits.
\end{corollary}

\subsection{Two scales of the determinant}

The same mathematical object --- $\det(G_h)$ for a triangle of 3 vertices --- operates at two scales:

\begin{enumerate}
\item \textbf{Network scale} (between simplices): the deficit angle $\delta_h = 2\pi - \sum\theta$ around a hinge shared by adjacent simplices. This measures curvature of the network.
\item \textbf{Internal scale} (within one simplex): the 10 hinges of a single simplex encode the internal structure. The AAA hinge, the 6 AAB hinges, and the 3 ABB hinges carry different geometric content.
\end{enumerate}

These are not separate structures. They are the same determinant evaluated at different levels of the simplicial hierarchy.

\subsection{Information flow between simplices}

All information between two adjacent simplices passes through their shared face. The face type determines which hinge types participate:

\begin{itemize}
\item Through an AAAB face: AAA $\times$ 1 + AAB $\times$ 3 = 4 bits. The AAA hinge participates.
\item Through an AABB face: AAB $\times$ 2 + ABB $\times$ 2 = 4 bits. No AAA hinge.
\end{itemize}

There is exactly one channel (the face) connecting adjacent simplices. All signals --- regardless of hinge type --- propagate through this single channel at the same rate. The ratio of B-direction to A-direction lattice densities is
\begin{equation}
c = \frac{n_B / d_B}{n_A / d_A} = \frac{2/1}{3/3} = 2,
\label{eq:c_geometric}
\end{equation}
where $d_A = n_A - (n_A - 1) = 1$ and $d_B = n_B - (n_B - 1) = 1$ are the respective codimensions within each sector. (The physical identification of $c$ with the speed of light is made in Chapter~\ref{ch:geometry}.)

\subsection{Vertex selection patterns}
\label{sec:fermion_combinatorics}

A \textbf{vertex selection} is a subset of the simplex's 5 vertices. Two types of selections are distinguished by their antisymmetry properties.

\subsubsection{1-vertex selections: the $\bar{5}$ pattern}

Selecting 1 vertex from 5: $\binom{5}{1} = 5$ patterns.

Each selected vertex participates in $\binom{4}{2} = 6$ of the 10 hinges (those containing the selected vertex). The hinge composition of each selection --- its \textbf{fingerprint} --- depends on the vertex type:

\begin{center}
\begin{tabular}{cccccc}
\hline
Selected & AAA & AAB & ABB & Total hinges & Count \\
\hline
A vertex & 1 & 4 & 1 & 6 & 3 \\
B vertex & 0 & 3 & 3 & 6 & 2 \\
\hline
\multicolumn{5}{r}{Total selections} & 5 \\
\hline
\end{tabular}
\end{center}

An A-type selection participates in the AAA hinge. A B-type selection does not.

\subsubsection{2-vertex selections: the $10$ pattern}

Selecting 2 vertices from 5: $\binom{5}{2} = 10$ patterns, antisymmetric under exchange ($\{v_i, v_j\} = -\{v_j, v_i\}$).

\begin{center}
\begin{tabular}{ccccc}
\hline
Pair type & Hinge fingerprint & Count & Involves AAA? \\
\hline
AA & AAA $\times$ 1 + AAB $\times$ 2 & $\binom{3}{2} = 3$ & Yes \\
AB & AAB $\times$ 2 + ABB $\times$ 1 & $3 \times 2 = 6$ & No \\
BB & ABB $\times$ 3 & $\binom{2}{2} = 1$ & No \\
\hline
\multicolumn{3}{r}{Total} & 10 \\
\hline
\end{tabular}
\end{center}

\subsubsection{Total: 15 antisymmetric patterns}

\begin{theorem}[Vertex selection count]
\label{thm:15}
The total number of antisymmetric vertex selections is
\begin{equation}
\binom{5}{1} + \binom{5}{2} = 5 + 10 = 15.
\end{equation}
\end{theorem}

The 1-vertex selection is odd-dimensional (a point) and inherits antisymmetric exchange statistics. The 2-vertex selection is antisymmetric by construction. The symmetric structures (edges, i.e., inner products $G_{ij}$) carry $\binom{5}{2} = 10$ symmetric patterns.

(The physical identification of these 15 antisymmetric patterns with the 15 Weyl fermions of one Standard Model generation, and the 10 symmetric patterns with gauge bosons, is made in Chapter~\ref{ch:geometry}.)

\subsection{The AAA closure condition}
\label{sec:confinement_geometry}

\begin{theorem}[AAA triangle closure]
\label{thm:aaa_closure}
The AAA hinge has 3 A-vertices. For $\det(G_{\text{AAA}}) > 0$, all 3 must be present. Removing any one vertex collapses the triangle: $\det \to 0$, destroying 1 bit of information.
\end{theorem}

\begin{proof}
A $3 \times 3$ Gram matrix of rank $\leq 2$ has $\det = 0$. Removing 1 of 3 vertices reduces the triangle to a line segment ($2 \times 2$ sub-matrix embedded in the original space), which has $\det_3 = 0$.
\end{proof}

\begin{corollary}
A-type vertices selected from the $\bar{5}$ and $10$ patterns that involve the AAA hinge are bound in triples. They cannot be isolated without destroying geometric information.
\end{corollary}

\begin{theorem}[BBB non-existence]
There is no analogous closure condition for B-type vertices: the BBB triangle does not exist ($\binom{2}{3} = 0$). B-type selections are free.
\end{theorem}

\begin{remark}
The number 3 (vertices of a triangle) equals $n_A$. This is not a coincidence: triangles are the codimension-2 objects in a space whose dimension is determined by the atomic pair $(n_A, n_B) = (3, 2)$.
\end{remark}

\subsection{Summary of geometric data}

\begin{center}
\begin{tabular}{rl}
\hline
Quantity & Value \\
\hline
Vertices per simplex & 5 \\
A-vertices & 3 \\
B-vertices & 2 \\
Faces per simplex & 5 (= 2 AAAB + 3 AABB) \\
Hinges per simplex & 10 (= 1 AAA + 6 AAB + 3 ABB) \\
Hinges per face & 4 \\
Faces per hinge & 2 \\
Bits per hinge & 1 \\
Bits per face & 4 \\
Bits per simplex & 10 \\
Antisymmetric selections & 15 (= 5 + 10) \\
Symmetric selections (edges) & 10 \\
Lattice density ratio $c$ & 2 \\
BBB hinges & 0 (impossible) \\
\hline
\end{tabular}
\end{center}

All numbers are determined by $d = 5$ and the atomic pair $(n_A, n_B) = (3, 2)$. No physical input was used.
